\documentclass[10pt]{article}
\usepackage{lingmacros}
\usepackage{tree-dvips}
\begin{document}
\pagenumbering{gobble}
\section*{Classificazione microrganismi}
\textbf{Saprofiti} come miceti che decompongono materiale organico rendendolo disponibile alle piante.\\
\textbf{Simbionti} come microrganismi che risiedono nel rumine dei ruminanti, gli azotofissatori delle leguminose e i licheni sulle alghe fotosintetiche/cianobatterio.\\
\textbf{Commensali} come microbiota intestinale e possono essere opportunisti (candida).\\
\textbf{Parassiti} come Colera (vibrione flagellato).\\
\textbf{Elminti} come vermi e tenie. \\
\section*{Malattie infettive}
L'\textbf{infezione} avviene principalmente sulle mucose del corpo (naso, bocca, ano, occhi).\\
L'\textbf{infestazione} è l'invasione in un organismo di parassiti e/o animali (pidocchi, elminti, scabbia).\\
\textbf{Dose Infettante 50\%} (DI50) è la quantità di microrganismi infettanti necessari a infettare il 50\% delle cavie. \\
\textbf{DL50} è la quantità di microrganismi infettanti necessari a uccidere il 50\% delle cavie. \\
\paragraph*{Malattia di Lyme} Causata da una spirocheta portata da una \underline{zecca} che succhia il sangue dai cervi o roditori, solamente le femmine sono ematofaghe perché aiuta con la produzione di uova. Provoca eritemi e può aggravare con il tempo con artrite e carditi fino a demielinizzazione neuronale con AD o sclerosi multipla. Grazie al loro pungiglione le zecche usano un movimento rotatorio quindi devono essere rimosse con un movimento contrario, inoltre si deve evitare i disinfettanti perché causa il rigurgito della zecca. \\
\paragraph*{Legionellosi} Causata dalla Legionella pneumophila (gram- opportunista) ha un'incidenza maggiore negli anziani e si riproduce negli \underline{ambienti acquatici} che risale attraverso tubature e aria condizionata dove forma \underline{biofilm}, non ha trasmissione persona per persona. La Legionella previene il killing intracellulare lidendo il fagolisosoma dopo essersi riprodotta all'interno. \\
\\
Un trattamento di antibiotici o chemio/radioterapici può causare una disbiosi intestinale portando a una rottura dell'equilibrio dell'ecosistema e favorendo lo sviluppo di candidosi. \\
Le peritoniti possono essere causate da un varco nell'intestino dovuto ad appendiciti che permette ai microrganismi intestinali a infiammare il peritoneo.\\
Le \textbf{zoonosi} possono fare un salto di specie diretto oppure usare un \underline{ospite intermedio} che entrando in contatto con un virus umano e di un'altra specie animale genera un virus ibrido. L'HIV è una zoonosi (spill-over) e nasce come virus delle scimmie.\\
La trasmissione di tubercolosi è dovuta all'espettorato e non semplice saliva, mentre altri microrganismi vengono trasmessi con \textbf{microgocce di Flugge}, le quali raggiungono dimensioni di 1-4 micron che gli permette di viaggiare in aria per molto tempo con al loro interno anche più di un solo microrganismo.\\
Il Treponema pallidum (sifilide) e la Neisseria gonorrhoeae (gonorrea) sono esclusivamente a trasmissione sessuale perché non sopravvivono facilmente all'esterno e usano il contatto tra mucose per trasmettersi. \\
L'infezione subclinica è un'infezione asintomatica come portatore sano. \\
\end{document}